\documentclass[11pt]{article}
\usepackage[margin=1in]{geometry}
\usepackage{amsmath, amssymb, amsthm}
\usepackage{booktabs}
\usepackage{enumitem}
\usepackage[hidelinks]{hyperref}

\newcommand{\cop}{\mathrm{cop}}
\newcommand{\pcf}{\mathrm{pcf}}
\newcommand{\Var}{\mathrm{Var}}
\newcommand{\diag}{\mathrm{diag}}
\newcommand{\code}[1]{\texttt{#1}}

\title{Why the conditional cells fail, and how to fix them}
\author{Diagnosis from the HANK validation exercise}
\date{September 12, 2026}

\begin{document}
\maketitle

\section{The problem in one sentence}

In the HANK exercise --- where the truth is known --- the model reproduces every
marginal statistic at 92--99\% correlation, and every conditional statistic that
is \emph{conditioned on} a well-observed measure at 92--99\%, but the cells that
report \textbf{wealth by income group} collapse to 0--15\% outside the bottom
half. Since the empirical application's headline numbers (relative consumption
and wealth of asset-rich, low-income households) are objects of exactly this
type, this has to be fixed rather than caveated.

The diagnosis below shows that the failure is \emph{not} a data limitation: the
information is present in the frequently observed series. It is a
representation failure concentrated in the copula block, with three identifiable
causes. Sections \ref{sec:notation}--\ref{sec:evidence} state the problem and
the evidence; Section \ref{sec:fixes} gives the fixes with their math (the
fourth is withdrawn on inspection of the code);
Section \ref{sec:order} gives the order of work.

\section{Notation}\label{sec:notation}

At each date $t$, dataset $j$ delivers a coefficient vector
\begin{equation}
  \theta_{j,t} \;=\; \bigl( \underbrace{\theta^{\cop}_{j,t}}_{n_c \approx 1694},\;
                            \underbrace{\theta^{\pcf}_{j,t}}_{n_q = D\cdot G = 36} \bigr)
  \;\in\; \mathbb{R}^{n}, \qquad n = n_c + n_q ,
\end{equation}
the Legendre coefficients of the copula density and of the $D$ marginal
percentile functions. The distributional block of the model is
\begin{align}
  \theta_{j,t} &= \mu_j + S_j\,\bigl( \Lambda f_t \bigr) + e_{j,t},
  &  e_{j,t} &\sim N(0,\Sigma_j), \label{eq:meas}\\
  \begin{pmatrix} f_t \\ a_t \end{pmatrix}
   &= \begin{pmatrix} A & B \\ C & D \end{pmatrix}
      \begin{pmatrix} f_{t-1} \\ a_{t-1} \end{pmatrix} + \eta_t,
  & \eta_t &\sim N(0,\Omega), \label{eq:state}
\end{align}
with $f_t \in \mathbb{R}^r$ the distributional factors ($r=5$ in the current
spec), $a_t \in \mathbb{R}^q$ the aggregate factors, $\Lambda \in
\mathbb{R}^{n\times r}$ the loadings (\code{proj}), and $\mu_j, S_j$ the
per-dataset demeaning and object-level scaling.

The statistics the paper reports are nonlinear functionals of $\theta$. In
discrete form, with $\pi_{jk}(t)$ the joint mass in (income group $j$, wealth
group $k$) implied by the copula and $w_k(t)$ the mean wealth of wealth group
$k$ implied by the percentile function,
\begin{equation}
  m_j(t) \;=\; \frac{\sum_k \pi_{jk}(t)\, w_k(t)}{\pi_{j\cdot}(t)},
  \qquad
  \bar W_t \;=\; \sum_{j,k} \pi_{jk}(t)\, w_k(t),
  \qquad
  \tilde m_j(t) \;=\; \frac{m_j(t)}{\bar W_t} .
  \label{eq:cell}
\end{equation}
$\tilde m_j$ is the plotted ``wealth relative to average, by income group''.

\section{Why these cells are the fragile ones}\label{sec:fragile}

Differentiating \eqref{eq:cell} splits the movement of a cell into two channels:
\begin{equation}
  \mathrm{d} m_j \;=\;
  \underbrace{\sum_k \frac{\pi_{jk}}{\pi_{j\cdot}}\, \mathrm{d}w_k}_{\text{marginal channel}}
  \;+\;
  \underbrace{\frac{1}{\pi_{j\cdot}}\sum_k \bigl( w_k - m_j \bigr)\, \mathrm{d}\pi_{jk}}_{\text{copula channel}} .
  \label{eq:decomp}
\end{equation}
Two consequences follow, and they explain the entire pattern in the results.

\paragraph{(i) Asymmetry.} The copula channel is weighted by $w_k - m_j$, the
dispersion of the \emph{conditioned} variable's group means. In the HANK economy
wealth group means span a factor of about $15$ across groups
($0.31 \to 4.5$ relative to average) while income group means span about $2$
($0.75 \to 1.6$). The same copula error $\mathrm{d}\pi$ is therefore amplified
roughly sevenfold in \emph{wealth-by-income} and attenuated in
\emph{income-by-wealth}. This is why one direction works and the other does not,
even though both use the same copula.

\paragraph{(ii) Which channel carries the time variation.} For
income-by-wealth, most movement over time comes from the income percentile
function moving, i.e.\ the marginal channel, which the model recovers at
92--99\%. A nearly static copula is therefore sufficient. For
wealth-by-income the wealth marginal moves slowly, so nearly all movement must
come from $\mathrm{d}\pi$ --- the copula channel. A static copula is fatal.

\paragraph{Aggregation.} Since $\sum_j \pi_{j\cdot} \tilde m_j = 1$ with weights
$(0.5, 0.4, 0.1)$, the bottom cell is tightly pinned while the top cell is
nearly free: an error of $2$ in the top-decile cell costs only $0.2$ spread over
the rest. Hence bottom-half cells are decent everywhere (52--82\%) and top cells
are uniformly bad (14--15\%).

\section{Evidence}\label{sec:evidence}

\paragraph{E1. The information is available.} Regressing each truth cell on
principal components of the truth's own marginal quantile grids --- the
information the model observes at high frequency --- gives:

\begin{center}
\begin{tabular}{lccc}
\toprule
truth cell & 5 PCs & $+1$ lag & 15 PCs \\
\midrule
\texttt{wealth\_by\_income\_bot50} & 0.96 & 0.97 & 0.99 \\
\texttt{wealth\_by\_income\_mid40} & 0.73 & 0.77 & 0.93 \\
\texttt{wealth\_by\_income\_top10} & 0.81 & 0.83 & 0.96 \\
\texttt{income\_by\_wealth} (all)  & 1.00 & 1.00 & 1.00 \\
\bottomrule
\end{tabular}
\end{center}

A five-factor model that spans the marginals therefore has enough information to
reach $R^2 \approx 0.75$--$0.85$ on the failing cells. It delivers $\approx 0$.

\paragraph{E2. The estimated copula is frozen.} From the reconstructed copula of
replica~a, the implied income--wealth rank correlation has mean $0.532$ with a
relative standard deviation of $2.3\%$, against $11.2\%$ for the truth's
dependence gradient, and the two paths correlate at only $0.32$. Freezing the
copula at its time average barely changes the reconstructed cell.

\paragraph{E3. The loadings come from twelve observations.} The pool passed to
the PCA is $P = [\,\theta_{j,t}\,]$ with columns indexed by (dataset, date), and
the code retains only columns with no missing entry,
\begin{equation}
  \mathcal{K} \;=\; \bigl\{ (j,t) : \text{all $n$ coefficients observed} \bigr\},
  \qquad
  \hat\Lambda \;=\; \text{top-$r$ left singular vectors of } P_{\mathcal{K}} .
\end{equation}
Only the dataset observing \emph{all} measures qualifies, at its wave dates:
$|\mathcal{K}| = 12$, hence $\mathrm{rank}(P_\mathcal{K}) \le 11$. Every
quarterly observation of a partially observed dataset is discarded at this step.
The empirical model has the identical structure (12 complete PSID columns,
1999--2021).

\paragraph{E4. The copula observations are then down-weighted to irrelevance.}
$\Sigma_j$ carries \emph{one} variance per object, broadcast across all
coefficients of that object:
\begin{equation}
  \Sigma_j \;=\; \diag\bigl( \sigma^2_{j,\cop}\mathbf{1}_{n_c},\;
                             \sigma^2_{j,1}\mathbf{1}_{G},\;\ldots,\;
                             \sigma^2_{j,D}\mathbf{1}_{G} \bigr).
  \label{eq:sigma}
\end{equation}
At the mode, the multiplier on replica~a's copula is $37.4$ against $1.3$--$12$
for every other object (prior mean $1$). The Kalman gain
$K_t = P_t G' (G P_t G' + \Sigma)^{-1}$ is then numerically zero on the copula
rows: the only three-dimensional copula observations in the system are ignored.

\paragraph{Why E4 is mechanical, and why it still matters.} With $r \ll n_c$ the
retained factors cannot span $\mathbb{R}^{n_c}$, so most copula coefficients are
unfittable \emph{by construction} and the single $\sigma^2_{j,\cop}$ in
\eqref{eq:sigma} absorbs the whole residual,
\begin{equation}
  \hat\sigma^2_{j,\cop} \;\approx\;
  \underbrace{\tfrac{1}{n_c}\bigl\| (I - P_\Lambda)\,\theta^{\cop}_{j} \bigr\|^2}_{\text{unspanned, large}}
  \;+\;
  \underbrace{\tfrac{1}{n_c}\bigl\| P_\Lambda\,\theta^{\cop}_{j} - \Lambda_{\cop} \hat f \bigr\|^2}_{\text{genuine fit error, small}},
  \qquad P_\Lambda = \Lambda_{\cop}(\Lambda_{\cop}'\Lambda_{\cop})^{-1}\Lambda_{\cop}' .
\end{equation}
The damage is that the informative directions are charged the same variance as
the uninformative ones: the components of the copula the factors \emph{can}
track are discarded along with the components they cannot.

\section{The four fixes}\label{sec:fixes}

\subsection{Fix 1: separate the spanned and orthogonal copula error}

\paragraph{Plain English.} Right now the model has a single ``how noisy is the
copula'' dial for roughly seventeen hundred numbers. Most of those numbers are
things five factors could never track, so the dial is turned up to ``ignore'' ---
and that throws away the handful of copula directions the factors \emph{can}
track. Give the copula two dials instead of one, or, equivalently, hand the
filter the few copula summaries it can actually use instead of all seventeen
hundred raw coefficients.

\paragraph{Math.} Replace \eqref{eq:sigma}'s copula block by
\begin{equation}
  \Var(e^{\cop}_{j,t}) \;=\; \sigma_{\parallel,j}^2\, P_\Lambda \;+\; \sigma_{\perp,j}^2\, (I - P_\Lambda),
  \label{eq:twosigma}
\end{equation}
one extra parameter per dataset. Equivalently and more cheaply, rotate the
observation: premultiply the copula rows by $\Lambda_{\cop}^{+} =
(\Lambda_{\cop}'\Lambda_{\cop})^{-1}\Lambda_{\cop}'$ and keep the resulting $r$
coordinates as the copula observables,
\begin{equation}
  \tilde y^{\cop}_{j,t} \;=\; \Lambda^{+}_{\cop}\, \theta^{\cop}_{j,t}
  \;=\; f_t + \Lambda^{+}_{\cop} e_{j,t},
  \qquad
  \Var(\Lambda^{+}_{\cop} e_{j,t}) = \sigma^2_{\parallel,j} (\Lambda_{\cop}'\Lambda_{\cop})^{-1},
\end{equation}
which reduces the copula block from $n_c$ rows to $r$ rows and removes the
unspanned variance from the likelihood entirely. The premultiplication machinery
already exists (\code{pre\_multiply\_part!}).

\paragraph{Expected effect.} The copula observations start moving the factors
again. This alone does not create copula dynamics the factors cannot represent,
so it is necessary but probably not sufficient.

\subsection{Fix 2: estimate the loadings from all blocks (two-step Bai--Ng, or EM)}

\paragraph{Plain English.} The loadings --- the map from factors to every
coefficient --- are currently fitted on twelve dates, because we throw away any
date where some measure is missing. But a CEX-type dataset tells us plenty about
consumption and income coefficients every quarter, and a SIPP-type dataset about
income and wealth. Use each dataset for what it does observe: first estimate the
factors from the large, always-observed block, then get each remaining
coefficient's loading by regressing it on those factors at whatever dates it
\emph{is} observed.

\paragraph{Math.} Let $\mathcal{B}$ index the balanced block (coefficients
observed at every $t$; in the replica exercise the quarterly consumption/income
rows) and $Y_\mathcal{B} \in \mathbb{R}^{N_B \times T}$ the corresponding panel.

\emph{Step 1 (factors).} With the normalization $F'F/T = I_r$,
\begin{equation}
  \hat F \;=\; \sqrt{T}\;\bigl[\text{eigenvectors of } Y_\mathcal{B}' Y_\mathcal{B}
              \text{ for the $r$ largest eigenvalues}\bigr],
  \qquad F = (f_1,\ldots,f_T)' .
\end{equation}

\emph{Step 2 (loadings for the sparse rows).} For each coefficient $i \notin
\mathcal{B}$ with observation set $\mathcal{T}_i$,
\begin{equation}
  \hat\lambda_i \;=\;
  \Bigl( \sum_{t \in \mathcal{T}_i} \hat f_t \hat f_t' \Bigr)^{-1}
  \sum_{t \in \mathcal{T}_i} \hat f_t\, \theta_{i,t} .
  \label{eq:step2}
\end{equation}
This is the standard two-pass estimator: $\hat\lambda_i$ is consistent as
$|\mathcal{T}_i| \to \infty$ with $\hat f_t$ consistent as $N_B, T \to \infty$.
Crucially, \eqref{eq:step2} is an $r$-regressor least-squares problem with
$|\mathcal{T}_i|$ observations, rather than a joint PCA on twelve columns.

\emph{ML alternative (EM / PPCA).} Iterate to a stationary point of the Gaussian
likelihood:
\begin{equation}
  \text{E-step: } \theta^{(s+1)}_{i,t} = \hat\lambda_i^{(s)\prime} \hat f^{(s)}_t
  \ \ \text{for } (i,t) \text{ missing};
  \qquad
  \text{M-step: } (\hat\Lambda^{(s+1)}, \hat F^{(s+1)}) = \mathrm{SVD}_r\bigl(\Theta^{(s+1)}\bigr).
\end{equation}

\paragraph{Caveat to state in the appendix.} The missingness is \emph{block}
structured: consumption and wealth coefficients co-occur only in the
PSID-type dataset. Cross-measure dependence information is therefore genuinely
limited to those dates, and no imputation scheme manufactures it. What Fix~2
recovers is the within-measure and two-dimensional-copula information from every
other dataset-period, which is currently discarded entirely. For the same
reason, pairwise-covariance PCA is not recommended: blocks that never co-occur
would be unidentified and the assembled matrix need not be positive
semi-definite.

\subsection{Fix 3: dedicated copula factors}

\paragraph{Plain English.} Even with better loadings, the copula has to share
five factors with the marginals, and the marginals win --- they are observed far
more often. Give the dependence structure a small number of factors of its own
that only the copula loads on, and let the data switch them on: if the copula
genuinely does not move, a shrinkage prior keeps them at zero and nothing
changes.

\paragraph{Math.} Partition $f_t = (f^m_t, f^c_t)$ with $f^c_t \in
\mathbb{R}^{r_c}$, $r_c$ small, and impose a block-triangular loading structure
in \eqref{eq:meas}:
\begin{equation}
  \begin{pmatrix} \theta^{\pcf}_{j,t} \\[2pt] \theta^{\cop}_{j,t} \end{pmatrix}
  \;=\;
  \begin{pmatrix} \Lambda_{qm} & 0 \\ \Lambda_{cm} & \Lambda_{cc} \end{pmatrix}
  \begin{pmatrix} f^m_t \\ f^c_t \end{pmatrix} + e_{j,t} ,
  \label{eq:blockload}
\end{equation}
so the copula-specific factors cannot contaminate the marginals. The zero block
also resolves the rotational indeterminacy between $f^m$ and $f^c$ up to
separate rotations within each block; normalize $\Lambda_{cc}'\Lambda_{cc}/n_c =
I_{r_c}$ as usual. The transition \eqref{eq:state} gains a block for $f^c$ with
its own persistence, which slots directly into the existing per-factor
persistence machinery ($\rho_{F,j}$ for $j$ in the copula block). Shrink the
copula factors' innovations with the same global--local prior used for the trend
states,
\begin{equation}
  \omega_{c,\ell} \sim N^{+}\bigl(0, (\lambda_c \psi_\ell)^2\bigr), \qquad
  \psi_\ell \sim C^{+}(0,1), \qquad
  \lambda_c \sim C^{+}(0, s_c),
\end{equation}
so that $r_c$ can be set generously without forcing the copula to move.

\paragraph{Why this is the one that should actually fix the cells.} By
\eqref{eq:decomp} the failing cells need $\mathrm{d}\pi \neq 0$ with the right
timing. Fixes 1 and 2 improve \emph{how well we measure} the copula; Fix 3 is
what gives the copula somewhere to move that is not already spoken for by the
marginals.

\subsection{Fix 4 (withdrawn): the reported cells are already grid-independent}

\paragraph{Plain English.} I originally proposed raising the HANK integration
grid from quintiles to deciles, on the theory that a top-decile group read off a
five-point grid must be interpolated. Checking the code, that is not what
happens: both the marginal group series and the cross-conditional cells are
computed by \emph{exact integration in coefficient space} over the true
$50/40/10$ boundaries, independently of the evaluation grid. The concern was
already addressed --- and there is a comment in the source recording that the
earlier ten-cell version was replaced precisely because its top cell froze the
within-decile composition that conditional wealth depends on.

\paragraph{What the code does.} The share-group series integrate the quantile
function over the exact boundaries,
\begin{equation}
  \texttt{share\_spec} = (0.5,\,0.4,\,0.1),
  \qquad
  \texttt{share\_intervals} = (0,\; 0.5,\; 0.9,\; 1),
\end{equation}
and the cross-conditional cells use the closed-form Legendre segment integrals
on the conditioning side together with a composite quadrature on the outcome
side,
\begin{equation}
  B_{s,k} = \int_{p_{s-1}}^{p_s} Q_k(v)\,\mathrm{d}v,
  \qquad
  A_j = \int_0^1 Q_j(u)\, V_t(u)\,\mathrm{d}u,
  \qquad
  m_s(t) = \frac{\sum_{j,k} \kappa_{jk}(t)\, A_j\, B_{s,k}}{\sum_k \kappa_{1k}(t)\, B_{s,k}},
\end{equation}
with $\kappa(t)$ the reconstructed copula coefficient tensor and $V_t$ the
outcome quantile function. Neither expression involves
\texttt{integral\_cop\_grid}. This also explains why both the truth cells and
the model cells satisfy the population-share identity
$0.5\,\tilde m_{\text{bot}} + 0.4\,\tilde m_{\text{mid}} + 0.1\,\tilde m_{\text{top}} \approx 1$.

\paragraph{What the grid does change.} (i) the quantile-level panels, whose
groups are grid-derived --- bottom40/next40/top20 at $G=5$, bottom50/next40/top10
at $G=10$; (ii) the exported functional-data grid ($G^D$ copula cells); and
(iii) the granularity of the interval/noise caches in the data layer
(\texttt{\_quintiles} vs \texttt{\_deciles}), which changes the
measurement-error calibration and therefore the estimates themselves. Because of
(iii), \textbf{likelihoods are not comparable across grids}, and raising the
grid introduces a confound into the experiment ladder without touching the
statistic under investigation. The HANK truth file also carries quintile
marginal series only, so the quantile-level panels cannot be validated at
$G=10$ without regenerating the truth.

\paragraph{Recommendation.} Keep $G=5$ for the copula experiments. Treat the
grid as a separate robustness question for the empirical model, where $G=10$ is
already the default.

\section{Order of work}\label{sec:order}

\begin{enumerate}[label=\textbf{\arabic*.}]
  \item \textbf{Fix 1} (contained: \code{SelectPrior.jl}, \code{Model.jl}).
        Cheap, and it immediately quantifies how much informative copula
        variance was being discarded: compare $\hat\sigma_\parallel$ with the
        current pooled $37.4$.
  \item \textbf{Fix 4 is withdrawn} --- the reported cells are computed by
        exact integration and do not depend on the grid; raising it would only
        change the data layer's noise calibration and confound the ladder.
  \item \textbf{Fix 2} (contained: the pool/PCA step in \code{ModelPrep.jl}).
        Biggest gain per unit of work, and it improves the empirical estimates
        for the same reason it improves the replica ones.
  \item \textbf{Fix 3} (genuine spec extension). Design note, implementation,
        adversarial review round, then a HANK run --- the same protocol used for
        the stochastic-trend block.
\end{enumerate}

\paragraph{Acceptance test.} After each step, re-run the replica exercise and
read the nine \texttt{wealth\_by\_income} cells. The target set by E1 is
$R^2 \approx 0.75$ on \texttt{mid40} and $\approx 0.8$ on \texttt{top10}; the
intermediate diagnostic is the relative standard deviation of the implied
income--wealth rank correlation, which must rise from $2.3\%$ toward the truth's
$11.2\%$. Marginal and income-by-wealth cells must not deteriorate: they are
currently at 92--99\% and are what the paper's other results rest on.

\end{document}
